% ===================================================================
% SKILL: KATEGORIK YARDIMCI BOX'LAR + KUTUSUZ ANA ANLATIM
% ===================================================================
% AMAC:
% Frame'i ana tasiyici olarak kullanmak; buyuk genel anlatimi kutu icine
% almamak; sadece not, dikkat ve sonuc gibi ikincil mesajlari renkli
% kutularla vurgulamak.
%
% NOT:
% Bu snippet yeni varsayimi tanimlar:
% 1. Ana konu anlatimi dogrudan frame govdesinde yazilir.
% 2. Box'lar yalnizca ikincil semantik vurgu icindir.
% 3. Cozum metni icin kirmizi yardimci komutlar saglanir.
% ===================================================================

% -------------------------------------------------------------------
% 1. PREAMBLE'A BIR KEZ EKLENECEK YARDIMCI TANIMLAR
% -------------------------------------------------------------------
\definecolor{probInfoGray}{HTML}{4B5563}
\definecolor{probResultOrange}{HTML}{FF8621}
\definecolor{probExampleGreen}{HTML}{2E7D32}

\setbeamercolor{block title example}{fg=white,bg=probExampleGreen}
\setbeamercolor{block body example}{fg=black,bg=probExampleGreen!8}

\newenvironment{infoblock}[1]
{
    \setbeamercolor{block title}{fg=white,bg=probInfoGray}
    \setbeamercolor{block body}{fg=black,bg=probInfoGray!10}
    \begin{block}{#1}
}
{
    \end{block}
}

\newenvironment{noteblock}[1]
{
    \begin{infoblock}{#1}
}
{
    \end{infoblock}
}

\newenvironment{resultblock}[1]
{
    \setbeamercolor{block title}{fg=white,bg=probResultOrange}
    \setbeamercolor{block body}{fg=black,bg=probResultOrange!12}
    \begin{block}{#1}
}
{
    \end{block}
}

\newcommand{\solutionghost}[1]{{\color{red!45!black}#1}}
\newcommand{\solutiontext}[1]{{\color{red!85!black}#1}}

% -------------------------------------------------------------------
% 2. HIZLI KURAL OZETI
% -------------------------------------------------------------------
\begin{frame}{Kullanim Mantigi}
    \small
    \begin{itemize}
        \item Frame zaten ana beyaz alandir; genel anlatimi buyuk bir block icine kapatma.
        \item Ana konu, tanim, teorem ve adim adim cozum dogrudan frame govdesinde yazilir.
        \item Box kullanilacaksa bunu sadece not, dikkat, sonuc veya kisa mikro ornek icin yap.
    \end{itemize}

    \begin{infoblock}{Hatirlatma}
        Bu gri kutu ana govdeyi tasimaz; yalnizca ek aciklama verir.
    \end{infoblock}
\end{frame}

% -------------------------------------------------------------------
% 3. KUTUSUZ KONU ANLATIMI + TEK DESTEK KUTUSU
% -------------------------------------------------------------------
\begin{frame}{Kutusuz Konu Anlatimi}
    \small
    Olasilikta once deneyi tanimlar, sonra ornek uzayi kurariz. Bir olay bu uzayin
    uygun bir alt kumesidir. Bu ana anlatim dogrudan frame govdesinde kalir.

    \vspace{0.3cm}
    \begin{itemize}
        \item Deney: fiziksel veya mantiksal surec
        \item Ornek uzay: tum olasi sonuc kumesi
        \item Olay: ornek uzayin ilgili alt kumesi
    \end{itemize}

    \begin{resultblock}{Ana Fikir}
        Kutu burada ana anlatimi tasimaz; sadece ana cikarimi ayri bir katmanda vurgular.
    \end{resultblock}
\end{frame}

% -------------------------------------------------------------------
% 4. COZUM METNI ICIN KIRMIZI AKIS
% -------------------------------------------------------------------
\begin{frame}{Cozum Renk Yardimcilari}
    \small
    Ilk cozum frame'inde dusuk kontrast:

    \[
        \solutionghost{P(A') = 1 - P(A)}
    \]

    Tam cozum frame'inde guclu kirmizi:

    \[
        \solutiontext{P(A') = 1 - P(A) = 1 - 0.12 = 0.88}
    \]
\end{frame}
